\begin{lemma}
Let $t\geq2$, let $K$ be a finite field of order $q\geq16$, where $q$ is a
power of two, and let $G=G_K(t)$ be the projective polarity graph.  Suppose
that $G$ has order $n$, degree $d$, and bilinear mixing parameter $\lambda$
with
\[
 |V(G)|=n,\qquad d_G(v)=d,\qquad \lambda\leq2\sqrt d,
\]
the bilinear spectral inequality, and
\[
 q^t/2\leq n\leq2q^t,\qquad q^{t-1}/2\leq d\leq2q^{t-1}.
\]
If
\[
20000t^2(t+1)\leq A,\quad 2tA\leq C,\quad4A\leq C,
\quad C\leq q,
\quad Cq(\operatorname{Nat.log}_2q)^2\leq k,
\]
then the improved pair digraph
\[
 V(D^*)=\{(a,b)\in V(G)^2:a\perp b\},\qquad
 ((a,b),(a',b'))\in A(D^*)
 \Longleftrightarrow a\perp b'\text{ and }a'\not\perp b
\]
has the following marked-tree certificate.  In the rooted tree whose level-$m$
vertices are forward-independent $m$-tuples in $D^*$, there is a marking
function $M$ assigning either marked or unmarked to every such tuple (the
root is marked) such that every vertex has at most $h=Aq^t$ marked children and
every root-to-vertex path has at most $w=Aq\operatorname{Nat.log}_2q$
unmarked vertices.  Thus the assignment is global over all levels of the
rooted tree, independently of the given $k$.
\end{lemma}
\begin{proof}
Write $L=\operatorname{Nat.log}_2q$.  Since $q$ is a power of two, say
$q=2^m$, the base-two logarithm is $L=m$.  The graph and the orthogonality
relation are those supplied by \noderef{PolarityGraph}.  Applying
\noderef{ExpanderMixing} to the bilinear inequality and indicator functions
gives, for all ordered subsets $X,Y\subseteq V(G)$,
\[
 \left|e_G(X,Y)-\frac dn|X||Y|\right|
 \leq\lambda\sqrt{|X||Y|}.                                  \tag{1}
\]
Loops are retained in this ordered edge count.

By \noderef{ForwardIndependentTuple} and
\noderef{ForwardIndependentCount}, a forward-independent tuple in $D^*$ is
the same as a consistent sequence
\[
 \sigma=(a_1,b_1,\ldots,a_m,b_m),\qquad
 a_i\perp b_i,\qquad
 a_i\perp b_j\Longrightarrow a_j\perp b_i\quad(i<j).       \tag{2}
\]
Indeed, the last implication is exactly the negation of the two conjuncts in
the arc rule.  Form the rooted tree of these consistent sequences, with the
empty sequence as root.

Fix a consistent sequence $\sigma$.  For each projective point $y$ define
\[
 W^\sigma(y)=\operatorname{span}\{b_i:a_i\perp y\},
 \qquad r^\sigma(y)=\dim W^\sigma(y).                       \tag{3}
\]
If $(a,b)$ extends $\sigma$, then $a\perp b$, and whenever $a_i\perp b$,
(2) gives $a\perp b_i$.  Thus $a\perp W^\sigma(b)$.  The orthogonal
complement has dimension $t+1-r^\sigma(b)$, so for $r^\sigma(b)\leq t$ the
number of possible $a$ for a fixed $b$ is at most
\[
 \frac{q^{t+1-r^\sigma(b)}-1}{q-1}\leq2q^{t-r^\sigma(b)},  \tag{4}
\]
and it is zero when $r^\sigma(b)=t+1$.

Put
\[
 U_r^\sigma=\{y:r^\sigma(y)\leq r\},\qquad
 Z_r^\sigma=U_r^\sigma\setminus U_{r-1}^\sigma,qquad U_{-1}^\sigma=\varnothing.
                                                                    \tag{5}
\]
For $0\leq r\leq t$ with $U_r^\sigma\ne\varnothing$, choose
$\ell=\ell_r^\sigma\leq r$ maximizing $|Z_\ell^\sigma|$.  Since the sets
$Z_0^\sigma,\ldots,Z_r^\sigma$ partition $U_r^\sigma$,
\[
 |Z_\ell^\sigma|\geq\frac{|U_r^\sigma|}{r+1}
 \geq\frac{|U_r^\sigma|}{2t}.                                \tag{6}
\]

Call $u\in Z_r^\sigma$ popular if it is contained in at least
$|Z_\ell^\sigma|/(16q)$ of the spaces $W^\sigma(y)$ with
$y\in Z_\ell^\sigma$.  Each such space has dimension $\ell$ and contains at
most $2q^{\ell-1}$ projective points.  Double-counting incidences $(u,y)$
therefore gives at most
\[
 \frac{|Z_\ell^\sigma|\,2q^{\ell-1}}
 {|Z_\ell^\sigma|/(16q)}=32q^\ell                       \tag{7}
\]
popular points.  Mark every child whose second coordinate is popular at its
level.  By (4), (7), and $\ell_r^\sigma\leq r$, this marks at most
\[
 \sum_{r=0}^t32q^{\ell_r^\sigma}2q^{t-r}
 \leq64(t+1)q^t\leq128tq^t.                                \tag{8}
\]

Call a first coordinate $a$ poor at level $r$ if, with
$\ell=\ell_r^\sigma$,
\[
 |N_G(a)\cap Z_\ell^\sigma|
 \leq |Z_\ell^\sigma|/(8q),                                \tag{9}
\]
and let $P_r$ be the poor set.  Then
\[
 e_G(P_r,Z_\ell^\sigma)
 \leq |P_r||Z_\ell^\sigma|/(8q).                             \tag{10}
\]
The hypotheses give $d/n\geq1/(4q)$, so every $a\in P_r$ satisfies
\[
 |N_G(a)\cap Z_\ell^\sigma|
 \leq d|Z_\ell^\sigma|/(2n).                                \tag{11}
\]
The positivity of $n,d$, (1), and the exact order and regularity hypotheses
are the hypotheses of the product estimate \noderef{AlonRodlBound}.  Applying
it to $P_r$ and $Z_\ell^\sigma$ gives
\[
 |P_r||Z_\ell^\sigma|
 \leq4\lambda^2n^2/d^2\leq1024q^{t+1}.                     \tag{12}
\]
Here we used $\lambda^2\leq4d$, $n\leq2q^t$, and
$d\geq q^{t-1}/2$; the direct calculation gives the stronger constant
$128q^{t+1}$.

Since $|Z_r^\sigma|\leq|Z_\ell^\sigma|$, another use of (1), together with
$d/n\leq4/q$ and $\lambda\leq4q^{(t-1)/2}$, yields
\[
\begin{aligned}
 e_G(P_r,Z_r^\sigma)
 &\leq\frac4q\,1024q^{t+1}
   +4q^{(t-1)/2}\sqrt{1024q^{t+1}}\\
 &\leq5000q^t.                                               \tag{13}
\end{aligned}
\]
Mark every child whose first coordinate is poor at the level
$r=r^\sigma(b)$.  Such a child has $a\perp b$ and $b\in Z_r^\sigma$, so it
is counted by $e_G(P_r,Z_r^\sigma)$.  Summing (13) over $r$ marks at most
$5000(t+1)q^t\leq10000tq^t$ children of this second type.

The two marking totals are at most $10128tq^t$.  The explicit assumption
$A\geq20000t^2(t+1)$, with $t\geq2$, implies
$10128t\leq A$.  Thus, marking the root and using both rules, every tree
vertex has at most
\[
                              h=Aq^t                         \tag{14}
\]
marked children.

It remains to prove the path bound.  Let $\sigma'$ be an unmarked extension
of $\sigma$ by $(a,b)$, put $r=r^\sigma(b)$ and
$\ell=\ell_r^\sigma$.  The child exists only when $r\leq t$.  Since $a$ is
not poor, at least $|Z_\ell^\sigma|/(8q)$ points $y\in Z_\ell^\sigma$
satisfy $a\perp y$.  Since $b$ is not popular, at most
$|Z_\ell^\sigma|/(16q)$ of those have $b\subseteq W^\sigma(y)$.  Hence at
least $|Z_\ell^\sigma|/(16q)$ points satisfy both conditions.  For each such
$y$, adjoining $b$ raises the dimension in (3) from $\ell$ to $\ell+1$.
Therefore
\[
 |U_\ell^{\sigma'}|
 \leq |U_\ell^\sigma|-\frac{|Z_\ell^\sigma|}{16q}
 \leq\left(1-\frac1{32tq}\right)|U_\ell^\sigma|.            \tag{15}
\]
Record this level $\ell$ on the unmarked extension.  The set being decreased
was nonempty immediately before the extension, because it contains
$Z_\ell^\sigma$.

Consider a root-to-vertex path and suppose it has more than
$w=AqL$ unmarked vertices.  Some $\ell\in[0,t]$ is recorded at least
$w/(t+1)$ times.  Along the path, each $U_\ell$ is nonincreasing, and (15)
applies at every recorded occurrence.  If $i^*$ is the last occurrence, then,
using $|U_\ell^{\sigma_0}|=n\leq2q^t$,
\[
 |U_\ell^{\sigma_{i^*-1}}|
 \leq2q^t\left(1-\frac1{32tq}\right)^{w/(t+1)-1}.              \tag{16}
\]
For completeness, $AqL/(t+1)\geq2$, so
$w/(t+1)-1\geq w/(2(t+1))$.  Since $1-x\leq e^{-x}$ and
$A\geq20000t^2(t+1)$, the factor in (16) is at most
\[
 \exp\left(-\frac{A L}{64t(t+1)}\right)
 \leq \exp(-300tL)<\frac1{4q^t},                            \tag{17}
\]
where $L\geq4$ and $q=2^L$ were used in the last inequality.  Thus (16) is
less than $1$, so the integer set $U_\ell^{\sigma_{i^*-1}}$ is empty.  This
contradicts the recorded nonemptiness immediately before the last extension.
Every root-to-vertex path therefore has at most
\[
                              w=AqL                         \tag{18}
\]
unmarked vertices.

The marking rule is defined for every consistent sequence, hence for every
vertex at every level of the rooted tree.  Mark the root, and mark precisely
the children selected by the popular-second-coordinate or poor-first-coordinate
rules above.  Equations (8) and (13) give the bound (14) on marked children at
each vertex, while the decay argument (15)--(17) gives (18) on the number of
unmarked vertices on every root-to-vertex path.  This is the claimed global
marking function; the Boolean signature count and the final estimate are
formed downstream from this certificate.
\end{proof}
